\chapter{METODE PENELITIAN}

\section{Dataset}

Penelitian ini menggunakan dataset primer yang disusun secara mandiri untuk merepresentasikan entitas dan aktivitas akademik pada domain Informatika dan Komputer (Infokom) di Universitas Negeri Surabaya (UNESA). Pengembangan dataset ini dilakukan untuk mengatasi keterbatasan ketersediaan basis data publik yang terintegrasi secara komprehensif. Perancangan atribut data difokuskan untuk mengoptimalisasi mekanisme \textit{Question-Answering} berbasis graf (\textit{GraphRAG}), yang mensyaratkan pelacakan asal-usul informasi (\textit{provenance tracking}) hingga ke dokumen asli guna memitigasi halusinasi \textit{LLM} \cite{zheng2024medsumgraph}.

Dataset dihimpun melalui integrasi data dari tiga sumber utama: (1) profil akademik dosen pada pangkalan data PDDIKTI \citep{pddikti_db}, (2) riwayat publikasi \textit{Google Scholar} \citep{google_scholar}, dan (3) metadata artikel ilmiah dari \textit{Scopus} \citep{scopus_db}. Data yang terkumpul mencakup informasi dari dosen di berbagai program studi di bawah rumpun Infokom UNESA, di antaranya: S1 Sains Data (13 dosen), S1 Kecerdasan Artifisial (5), S1 Teknik Informatika (9), S1 Sistem Informasi (10), S2 Informatika (6), dan S1 Teknik Elektro (24). Selain itu, dataset ini juga mencakup data dari program studi terkait seperti S1 Bisnis Digital dan kependidikan teknik lainnya untuk memperkaya cakupan \textit{knowledge base}. Rincian struktur atribut disajikan pada Tabel \ref{tab:deskripsi-dataset}.

\begin{longtable}{|P{0.15\linewidth}|P{0.22\linewidth}|P{0.13\linewidth}|P{0.4\linewidth}|}
    \caption{Struktur Atribut Dataset Penelitian} \label{tab:deskripsi-dataset}                                                      \\

    \hline
    \textbf{Dataset} & \textbf{Nama Atribut} & \textbf{Tipe Data} & \textbf{Deskripsi}                                               \\
    \hline
    \endfirsthead

    \multicolumn{4}{c}{{\textit{\ref{tab:deskripsi-dataset} (lanjutan).}}}                                                           \\
    \hline
    \textbf{Dataset} & \textbf{Nama Atribut} & \textbf{Tipe Data} & \textbf{Deskripsi}                                               \\
    \hline
    \endhead

    \hline
    \multicolumn{4}{r}{\textit{Bersambung ke halaman berikutnya}}                                                                    \\
    \endfoot
    \hline
    \endlastfoot

    \multirow{9}{2.2cm}{\centering \textbf{Data Profil Dosen}}
                     & nama\_dosen           & String             & Nama lengkap dosen beserta gelar terverifikasi.                  \\* \hhline{|~|---|}
                     & nama\_norm            & String             & Nama dosen yang dinormalisasi tanpa gelar.                       \\* \hhline{|~|---|}
                     & nip                   & String             & Nomor Induk Pegawai.                                             \\* \hhline{|~|---|}
                     & nidn                  & String             & Nomor Induk Dosen Nasional.                                      \\* \hhline{|~|---|}
                     & prodi                 & String             & Program studi penempatan.                                        \\* \hhline{|~|---|}
                     & scholar\_id           & String             & ID unik profil Google Scholar dosen.                             \\* \hhline{|~|---|}
                     & scopus\_id            & String             & ID unik profil Scopus dosen.                                     \\* \hhline{|~|---|}
                     & sinta\_id             & String             & ID unik profil SINTA dosen.                                      \\* \hhline{|~|---|}
                     & source                & String             & Sumber akuisisi data (mis.\ \textit{WEB+PDDIKTI}, \textit{WEB}). \\
    \hline \pagebreak \hline

    \multirow{11}{2.2cm}{\centering \textbf{Data}                                                                                    \\ \textbf{Publikasi} \\ \textbf{Ilmiah} \\ \footnotesize \textit{(Scopus \& Scholar)}}
                     & Authors               & String             & Daftar nama penulis artikel ilmiah.                              \\* \hhline{|~|---|}
                     & Author IDs            & String             & Kumpulan ID profil penulis dari paper.                           \\* \hhline{|~|---|}
                     & Title                 & String             & Judul artikel ilmiah publikasi.                                  \\* \hhline{|~|---|}
                     & Year                  & Integer            & Tahun publikasi untuk kebutuhan pencarian temporal.              \\* \hhline{|~|---|}
                     & Journal               & String             & Nama jurnal, konferensi, atau prosiding publikasi.               \\* \hhline{|~|---|}
                     & Link                  & String             & Tautan langsung menuju halaman \textit{publisher} atau PDF.      \\* \hhline{|~|---|}
                     & Abstract              & Text               & Ringkasan isi artikel untuk rujukan utama RAG.                   \\* \hhline{|~|---|}
                     & Keywords              & String             & Kata kunci asli dari penulis untuk perkuatan \textit{graph}.     \\* \hhline{|~|---|}
                     & Document Type         & String             & Jenis publikasi (mis.\ \textit{Article}, \textit{Conference}).   \\* \hhline{|~|---|}
                     & DOI                   & String             & \textit{Digital Object Identifier} sebagai resolusi identitas.   \\* \hhline{|~|---|}
                     & TLDR                  & Text               & \textit{Too Long; Didn't Read}, intisari singkat oleh AI.        \\
\end{longtable}

\section{Alur Penelitian}

Alur penelitian dirancang sebagai rangkaian proses sistematis untuk membangun \textit{Knowledge Graph} dan mengimplementasikan \textit{Hybrid Vector-Graph Retrieval} dalam pencarian literatur domain Informatika dan Komputer di UNESA. Gambaran umum alur ini ditunjukkan pada \ref{fig:flowchart-utama}.

% ============================================
% FLOWCHART METODOLOGI TUGAS AKHIR
% Knowledge Graph-RAG untuk Penelusuran Literatur Akademik UNESA
% ============================================

% Definisi Style Flowchart (ANSI Standard)
% - Terminator  : Rounded rectangle (Mulai/Selesai)
% - Process     : Rectangle (aktivitas umum)
% - Predefined  : Rectangle + garis vertikal (sub-proses berdetail)
% - Decision    : Diamond (percabangan Ya/Tidak)
% - Off-page    : Pentagon (rujukan ke diagram detail)
% - Arrow       : Solid = alur utama, Dashed = rujukan detail
\tikzset{
    startstop/.style={rectangle, rounded corners=5pt, minimum width=2.8cm, minimum height=0.8cm, text centered, draw=black, line width=0.6pt, fill=green!20, font=\scriptsize\bfseries},
    process/.style={rectangle, minimum width=2.8cm, minimum height=0.8cm, text centered, draw=black, line width=0.6pt, fill=blue!12, font=\scriptsize},
    predefined/.style={rectangle, minimum width=2.8cm, minimum height=0.8cm, text centered, draw=black, line width=0.6pt, fill=orange!20, font=\scriptsize,
            path picture={\draw[line width=0.5pt] ([xshift=3pt]path picture bounding box.north west) -- ([xshift=3pt]path picture bounding box.south west);
                    \draw[line width=0.5pt] ([xshift=-3pt]path picture bounding box.north east) -- ([xshift=-3pt]path picture bounding box.south east);}},
    data/.style={trapezium, trapezium left angle=70, trapezium right angle=110, minimum width=1.8cm, minimum height=0.6cm, text centered, draw=black, line width=0.6pt, fill=yellow!20, font=\scriptsize},
    database/.style={cylinder, shape border rotate=90, draw=black, line width=0.6pt, minimum height=0.8cm, minimum width=1.2cm, fill=purple!15, font=\scriptsize, aspect=0.35},
    decision/.style={diamond, minimum width=1.8cm, minimum height=1.2cm, text centered, draw=black, line width=0.6pt, fill=red!12, font=\scriptsize, inner sep=1pt, aspect=2},
    onpage/.style={circle, minimum size=0.5cm, draw=black, line width=0.6pt, fill=white, font=\scriptsize\bfseries},
    offpage/.style={shape=signal, signal to=south, draw, fill=white, drop shadow, text=black, font=\bfseries\scriptsize, inner sep=3pt, minimum width=0.5cm, minimum height=0.5cm},
    arrow/.style={thick, -Stealth, line width=0.5pt},
    subbox/.style={draw=orange!70!black, line width=0.8pt, rounded corners=3pt, fill=orange!5, inner sep=8pt}
}

% FLOWCHART UTAMA — Vertikal Top-to-Bottom (ANSI Standard)
\begin{figure}[htbp]
    \centering
    \begin{tikzpicture}[node distance=0.7cm, scale=1, transform shape]

        % === ALUR UTAMA (Vertikal) ===
        \node (start) [startstop] {Mulai};

        \node (datacoll) [predefined, below=of start, text width=2.4cm] {Pengumpulan Data};

        \node (dec1) [decision, below=of datacoll, text width=1.5cm] {Data\\Lengkap?};

        \node (dataprep) [predefined, below=of dec1, text width=2.4cm] {Persiapan Data};

        \node (kgconst) [predefined, below=of dataprep, text width=2.4cm] {Konstruksi KG};

        \node (ragsys) [predefined, below=of kgconst, text width=2.4cm] {Hybrid\\GraphRAG};

        \node (eval) [process, below=of ragsys, text width=2.4cm] {Evaluasi Sistem};

        \node (dec2) [decision, below=of eval, text width=1.5cm] {Memenuhi\\Threshold?};

        \node (gui) [process, below=of dec2, text width=2.4cm] {Implementasi GUI};

        \node (endnode) [startstop, below=of gui] {Selesai};

        % === OFF-PAGE CONNECTORS (konsisten di kanan) ===
        \node (refA) [offpage, right=1.0cm of datacoll] {A};
        \node (refB) [offpage, right=1.0cm of dataprep] {B};
        \node (refC) [offpage, right=1.0cm of kgconst] {C};
        \node (refD) [offpage, right=1.0cm of ragsys] {D};

        % Garis putus-putus: proses → connector (rujukan ke diagram detail)
        \draw [densely dashed, ->] (datacoll) -- (refA);
        \draw [densely dashed, ->] (dataprep) -- (refB);
        \draw [densely dashed, ->] (kgconst) -- (refC);
        \draw [densely dashed, ->] (ragsys) -- (refD);

        % === ALUR UTAMA (Panah solid ke bawah) ===
        \draw [arrow] (start) -- (datacoll);
        \draw [arrow] (datacoll) -- (dec1);
        \draw [arrow] (dec1) -- node[right, font=\scriptsize] {Ya} (dataprep);
        \draw [arrow] (dataprep) -- (kgconst);
        \draw [arrow] (kgconst) -- (ragsys);
        \draw [arrow] (ragsys) -- (eval);
        \draw [arrow] (eval) -- (dec2);
        \draw [arrow] (dec2) -- node[right, font=\scriptsize] {Ya} (gui);
        \draw [arrow] (gui) -- (endnode);

        % === LOOP 1: Data tidak lengkap → kembali ke Pengumpulan Data ===
        \draw [arrow] (dec1.west) -- node[above, font=\scriptsize] {Tidak} ++(-1.5,0)
        |- (datacoll.west);

        % === LOOP 2: Evaluasi tidak memenuhi → kembali ke Hybrid GraphRAG ===
        \draw [arrow] (dec2.west) -- node[above, font=\scriptsize] {Tidak} ++(-1.5,0)
        |- (ragsys.west);

    \end{tikzpicture}
    \caption{Alur metodologi penelitian}
    \label{fig:flowchart-utama}
\end{figure}

\begin{enumerate}[label=\textbf{\arabic*.}, leftmargin=*, itemsep=10pt]

    % 1. Data Collection
    \item \textbf{Pengumpulan Data (\textit{Data Collection})}

          Tahap ini menghimpun dua kategori data utama: (1) data profil dosen beserta identitas indeksasinya, dan (2) data publikasi ilmiah beserta metadata pendukung. Kedua kategori ini menjadi fondasi untuk membangun \textit{Scholarly Knowledge Graph} (\textit{SKG}) yang merepresentasikan ekosistem publikasi Infokom UNESA. Proses akuisisi data dilakukan secara bertahap melalui \textit{pipeline} modular yang mengintegrasikan enam sumber data, sebagaimana ditunjukkan pada \ref{fig:detail-data-collection}.

          \begin{figure}[htbp]
              \centering
              \begin{tikzpicture}[scale=1, transform shape]
                  \node [subbox, fit={(-0.2,0.2) (4.8,-6.3)}, label={[font=\scriptsize\bfseries]above:Detail Data Collection}] {};

                  \node (subA) [offpage] at (2.3, -0.1) {A};
                  % Baris 1: Sumber Data Profil
                  \node (db1) [database, font=\tiny] at (0.6, -1.3) {Web Prodi};
                  \node (db2) [database, font=\tiny] at (2.3, -1.3) {PDDikti};
                  \node (db3) [database, font=\tiny] at (4.0, -1.3) {SimCV};
                  % Baris 2: Sumber ID \& Publikasi
                  \node (db4) [database, font=\tiny] at (0.6, -2.5) {SINTA};
                  \node (db5) [database, font=\tiny] at (2.3, -2.5) {Scholar};
                  \node (db6) [database, font=\tiny] at (4.0, -2.5) {Scopus};
                  % Pipeline
                  \node (pipe) [process, text width=2.2cm] at (2.3, -3.7) {Pipeline\\Modular};
                  \node (rawdata) [data] at (2.3, -4.9) {Raw Dataset};
                  \node (subB) [offpage] at (2.3, -6.0) {B};

                  \draw [arrow] (subA) -- (db2);
                  \draw [arrow] (db1) -- (pipe);
                  \draw [arrow] (db2) -- (pipe);
                  \draw [arrow] (db3) -- (pipe);
                  \draw [arrow] (db4) -- (pipe);
                  \draw [arrow] (db5) -- (pipe);
                  \draw [arrow] (db6) -- (pipe);
                  \draw [arrow] (pipe) -- (rawdata);
                  \draw [arrow] (rawdata) -- (subB);
              \end{tikzpicture}
              \caption{Detail proses \textit{Data Collection}}
              \label{fig:detail-data-collection}
          \end{figure}

          \begin{enumerate}[label=\textbf{\alph*.}, leftmargin=0.75cm, itemsep=5pt]
              \item \textbf{Data Profil Dosen}

                    Profil dosen dikumpulkan melalui \textit{pipeline} bertahap yang menggabungkan beberapa sumber. Laman web resmi setiap program studi UNESA dijadikan sebagai sumber utama (\textit{source of truth}) karena memuat nama lengkap beserta gelar, tautan profil akademik (\textit{Google Scholar}, \textit{Scopus}, \textit{SINTA}), serta NIP dosen. Data ini kemudian diperkaya secara bertahap: API PDDikti \citep{pddikti_db} menyediakan NIDN dan NIP yang otoritatif, dilanjutkan dengan verifikasi identitas melalui \textit{SimCV UNESA} (\textit{cv.unesa.ac.id}) untuk memastikan keakuratan nama dan afiliasi program studi.

                    Selanjutnya, ID indeksasi dosen yang belum tersedia dari sumber sebelumnya dilengkapi melalui halaman departemen SINTA \citep{sinta_database} untuk memperoleh \textit{SINTA ID}, \textit{SciVal} \citep{scopus_db} untuk \textit{Scopus ID}, serta verifikasi profil \textit{Google Scholar} \citep{google_scholar} untuk \textit{Scholar ID}. Setiap tahap melakukan pencocokan nama (\textit{fuzzy matching}) untuk mengasosiasikan data dari sumber berbeda ke entitas dosen yang sama, sehingga menghindari duplikasi.

              \item \textbf{Data Publikasi Ilmiah}

                    Data publikasi dihimpun dari dua sumber utama. Pertama, metadata publikasi terindeks \textit{Scopus} diekspor secara otomatis melalui antarmuka \textit{SciVal} \citep{scopus_db} untuk setiap dosen yang telah memiliki \textit{Scopus ID}. Data yang diperoleh meliputi judul, penulis, tahun, jurnal, DOI, dan tipe dokumen.

                    Kedua, publikasi dari \textit{Google Scholar} \citep{google_scholar} diakuisisi menggunakan \textit{SerpAPI} berdasarkan \textit{Scholar ID} masing-masing dosen. Sebelum disimpan, dilakukan deduplikasi silang terhadap data \textit{Scopus} untuk mencegah data ganda. Seluruh publikasi kemudian diperkaya metadatanya: \textit{CrossRef API} digunakan untuk memverifikasi DOI dan memperoleh kata kunci, \textit{scraping} halaman \textit{publisher} untuk mengekstrak abstrak, serta \textit{Semantic Scholar API} untuk menghasilkan ringkasan singkat (\textit{TLDR}). Strategi pengayaan bertingkat ini memaksimalkan kelengkapan data meskipun tidak semua sumber memberikan informasi yang sama.
          \end{enumerate}

          % 2. Data preparation
    \item \textbf{Persiapan Data (\textit{Data Preparation})}

          \hspace*{0.75cm}Data mentah dari tahap sebelumnya selanjutnya diproses untuk menjamin kualitas \textit{retrieval} (\ref{fig:detail-data-preparation}), melalui tahapan berikut:

          \begin{figure}[htbp]
              \centering
              \begin{tikzpicture}[scale=1, transform shape]
                  \node [subbox, fit={(0,0.2) (3.8,-6.3)}, label={[font=\scriptsize\bfseries]above:Detail Data Preparation}] {};

                  \node (subB2) [offpage] at (1.9, -0.1) {B};
                  \node (std1) [process, text width=2.2cm] at (1.9, -1.3) {Data Cleaning\\\& Normalization};
                  \node (std2) [process, text width=2.2cm] at (1.9, -2.5) {Text\\Chunking};
                  \node (std3) [process, text width=2.2cm] at (1.9, -3.7) {Embedding\\Generation};
                  \node (cleandata) [data, text width=1.8cm] at (1.9, -4.9) {Clean Data};
                  \node (subOff) [offpage] at (1.9, -6.0) {C};

                  \draw [arrow] (subB2) -- (std1);
                  \draw [arrow] (std1) -- (std2);
                  \draw [arrow] (std2) -- (std3);
                  \draw [arrow] (std3) -- (cleandata);
                  \draw [arrow] (cleandata) -- (subOff);
              \end{tikzpicture}
              \caption{Detail proses \textit{Data Preparation}}
              \label{fig:detail-data-preparation}
          \end{figure}

          \begin{enumerate}[label=\textbf{\alph*.}, leftmargin=0.75cm, itemsep=4pt]
              \item \textbf{\textit{Data Cleaning} \& \textit{Normalization}:} Menangani \textit{missing values} dan menstandarisasi variasi penulisan nama institusi serta singkatan gelar untuk mencegah duplikasi entitas dalam graf.

              \item \textbf{\textit{Text Chunking} (Pemotongan Teks):} Dokumen panjang seperti abstrak dipecah menjadi segmen semantik (misalnya 512 token) menggunakan \textit{semantic chunking}, mengingat keterbatasan \textit{context window} \textit{LLM}. Strategi ini menjaga keutuhan unit informasi sekaligus mengurangi \textit{noise}.

              \item \textbf{\textit{Embedding Generation}:} Setiap \textit{chunk} dikonversi menjadi vektor berdimensi tinggi menggunakan model \textit{embedding} (misalnya \textit{all-MiniLM-L6-v2}), memungkinkan pencarian berbasis makna (\textit{semantic search}) yang melampaui pencocokan kata kunci leksikal.
          \end{enumerate}


          % 3. Construction of the Knowledge Graph
    \item \textbf{Konstruksi \textit{Knowledge Graph}}

          \hspace*{0.75cm}Tahap ini mentransformasi data yang telah dibersihkan menjadi basis pengetahuan terstruktur untuk penalaran relasional. Seperti ditunjukkan pada \ref{fig:detail-kg-construction}, data terstruktur (metadata dosen/publikasi) dan tidak terstruktur (abstrak) diproses melalui dua jalur paralel yang disatukan melalui \textit{Entity Linking}.

          \begin{figure}[htbp]
              \centering
              \begin{tikzpicture}[scale=1, transform shape]
                  % Outer box — widened to fit all 2.8cm-wide process nodes with margin
                  \node [subbox, fit={(-1.0,0.3) (8.0,-12.9)}, label={[font=\scriptsize\bfseries]above:Detail KG Construction}] {};

                  \node (subC2) [offpage] at (3.0, -0.1) {C};

                  % 1. Ontology Modeling
                  \node (ontmod) [process, text width=2.2cm] at (3.0, -1.3) {Ontology\\Modeling};

                  \node (input) [process, text width=2.2cm] at (3.0, -2.6) {Parallel\\Processing};

                  % === TWO INPUT SOURCES ===
                  \node (structured) [data, fill=gray!30, text width=1.4cm] at (1.0, -3.9) {Structured\\Data};
                  \node (unstructured) [data, fill=gray!30, text width=1.4cm] at (5.5, -3.9) {Unstructured\\Data};

                  % === STRUCTURED PATH (Left) ===
                  \node (entres_struct) [process, text width=1.8cm] at (1.0, -5.3) {Entity\\Resolution};

                  % === UNSTRUCTURED PATH (Right) ===
                  \node (textproc) [process, text width=1.8cm] at (5.5, -5.3) {Text\\Processing};
                  \node (coref) [process, text width=1.8cm] at (5.5, -6.7) {Coreference\\Resolution};

                  % NER & RE — shifted outward, explicit min-width to avoid 2.8cm default
                  \node (ner) [process, minimum width=1.9cm, minimum height=0.7cm, text width=1.1cm] at (4.5, -8.1) {NER\\(GLiNER)};
                  \node (re) [process, minimum width=1.9cm, minimum height=0.7cm, text width=1.1cm] at (6.5, -8.1) {RE\\(GLiRel)};
                  \node (lexgraph) [data, text width=1.6cm] at (5.5, -9.4) {Lexical\\Graph};

                  % Entity Linking — moved left for clear separation from Lexical Graph
                  \node (entlink) [process, text width=2.2cm] at (2.0, -9.4) {Entity Linking};

                  % Indexing
                  \node (indexing) [process, text width=2.2cm] at (3.0, -10.7) {Dual Indexing};

                  % === OUTPUT DATABASES ===
                  \node (graphdb) [database] at (1.5, -11.9) {Graph DB};
                  \node (vectordb) [database] at (4.5, -11.9) {Vector DB};
                  \node (subD2) [offpage] at (3.0, -12.6) {D};

                  % === ARROWS - Input Split ===
                  \draw [arrow] (subC2) -- (ontmod);
                  \draw [arrow] (ontmod) -- (input);
                  \draw [arrow] (input) -| (structured);
                  \draw [arrow] (input) -| (unstructured);

                  % === ARROWS - Structured Path (Left) ===
                  \draw [arrow] (structured) -- (entres_struct);
                  \draw [arrow] (entres_struct.south) -- (1.0, -7.5) -| (entlink.north);

                  % === ARROWS - Unstructured Path (Right) ===
                  \draw [arrow] (unstructured) -- (textproc);
                  \draw [arrow] (textproc) -- (coref);
                  \draw [arrow] (coref) -- (ner);
                  \draw [arrow] (coref) -- (re);
                  \draw [arrow] (ner) -- (lexgraph);
                  \draw [arrow] (re) -- (lexgraph);
                  \draw [arrow] (lexgraph) -- (entlink);
                  \draw [arrow] (entlink.south) -- (2.0, -10.1) -| (indexing.north);
                  \draw [arrow] (indexing) -- (graphdb);
                  \draw [arrow] (indexing) -- (vectordb);
                  \draw [arrow] (graphdb) |- (subD2);
                  \draw [arrow] (vectordb) |- (subD2);
              \end{tikzpicture}
              \caption{Detail tahapan konstruksi \textit{Knowledge Graph}}
              \label{fig:detail-kg-construction}
          \end{figure}

          \begin{enumerate}[label=\textbf{\arabic*.}, leftmargin=0.75cm, itemsep=5pt]
              \item \textbf{\textit{Ontology Modeling (Pemodelan Skema)}}

                    \begin{figure}[htbp]
                        \centering
                        \begin{tikzpicture}[
                                >=Stealth,
                                scale=1, transform shape,
                                % Node Styles with precise colors and shapes
                                base/.style={draw=gray!50, line width=0.8pt, align=center, inner sep=2pt, font=\large\sffamily},
                                % Metadata Nodes (Rectangles) - Consistent Size (2.2cm x 1.2cm)
                                titlenode/.style={base, rectangle, rounded corners=6pt, fill=teal!35, draw=teal!60!black, minimum width=2.2cm, minimum height=1.2cm, font=\large\bfseries},
                                authornode/.style={base, rectangle, rounded corners=6pt, fill=violet!15, draw=violet!50!black, minimum width=2.2cm, minimum height=1.2cm},
                                instnode/.style={base, rectangle, rounded corners=6pt, fill=blue!15, draw=blue!50!black, minimum width=2.2cm, minimum height=1.2cm},
                                datenode/.style={base, rectangle, rounded corners=6pt, fill=cyan!30, draw=cyan!60!black, minimum width=2.2cm, minimum height=1.2cm},
                                confnode/.style={base, rectangle, rounded corners=6pt, fill=lime!30, draw=lime!60!black, minimum width=2.2cm, minimum height=1.2cm},
                                keynode/.style={base, rectangle, rounded corners=5pt, fill=green!10, draw=green!50!black, minimum width=2.2cm, minimum height=1.2cm},
                                fieldnode/.style={base, rectangle, rounded corners=6pt, fill=cyan!20, draw=cyan!50!black, minimum width=2.2cm, minimum height=1.2cm},
                                % Scientific Concepts (Circles) - Consistent Size (1.8cm)
                                problemnode/.style={base, circle, fill=orange!30, draw=orange!60!black, minimum size=1.8cm},
                                methodnode/.style={base, circle, fill=yellow!40, draw=yellow!60!black, minimum size=1.8cm},
                                metricnode/.style={base, circle, fill=magenta!20, draw=magenta!60!black, minimum size=1.8cm},
                                datasetnode/.style={base, circle, fill=violet!20, draw=violet!60!black, minimum size=1.8cm},
                                tasknode/.style={base, circle, fill=red!30, draw=red!60!black, minimum size=1.8cm},
                                modelnode/.style={base, circle, fill=red!15, draw=red!50!black, minimum size=1.8cm},
                                resultsnode/.style={base, circle, fill=blue!10, draw=blue!50!black, minimum size=1.8cm},
                                innovnode/.style={base, circle, fill=blue!25, draw=blue!60!black, minimum size=1.8cm},
                                % Connection Styles
                                arrow/.style={->, line width=0.7pt, gray!60!black},
                                labelstyle/.style={font=\footnotesize\sffamily, fill=white, text=gray!40!black, align=center, inner sep=1pt, fill opacity=0.9, text opacity=1, rounded corners=1pt}
                            ]

                            % ==========================================================
                            % NODE PLACEMENT (Compact & Consistent)
                            % ==========================================================

                            % Center
                            \node[titlenode] (Title) at (0, 0) {Title};

                            % Top Arc (Metadata)
                            \node[datenode] (Date) at (-5.0, 0.0) {Date};
                            \node[authornode] (Author) at (-3.9, 1.5) {Author};
                            \node[instnode] (Inst) at (-1, 2.8) {Institution};
                            \node[confnode] (Conf) at (1.7, 2.8) {Conference};
                            \node[keynode] (Key) at (4, 1.5) {Keywords};
                            \node[fieldnode] (Field) at (5, 0.0) {Field};

                            % Middle Row (Concepts) - Compacted Layout
                            \node[problemnode] (Prob) at (-6.0, -2.5) {Problem};
                            \node[methodnode] (Method) at (-2.5, -2.5) {Method};
                            \node[metricnode] (Metric) at (3.0, -2.5) {Metric};
                            \node[datasetnode] (Dataset) at (6.0, -2.5) {Dataset};

                            % Bottom Row (Outputs) - Compacted & Aligned for Straight Lines
                            \node[tasknode] (Task) at (-4.5, -6.0) {Task};
                            \node[modelnode] (Model) at (-1.5, -6.0) {Model};
                            \node[resultsnode] (Results) at (1.5, -6.0) {Results};
                            \node[innovnode] (Innov) at (4.5, -6.0) {Innovation};

                            % ==========================================================
                            % CONNECTIONS (Straight Lines)
                            % ==========================================================

                            % --- Connections that need to be behind nodes ---
                            \begin{scope}[on background layer]
                                \draw[arrow] (Title) -- node[labelstyle, sloped, near end] {works on} (Task);
                                \draw[arrow] (Title) -- node[labelstyle, sloped, pos=0.25] {innovates} (Innov);
                            \end{scope}

                            % --- Metadata Connections ---
                            \draw[arrow] (Author.north) -- node[labelstyle, sloped, pos=0.5] {works for} (Inst.west);
                            \draw[arrow] (Author) -- node[labelstyle, sloped] {writes} (Title);
                            \draw[arrow] (Title) -- node[labelstyle, sloped] {writes on} (Date);
                            \draw[arrow] (Conf) -- node[labelstyle, sloped] {publishes} (Title);
                            \draw[arrow] (Title) -- node[labelstyle, sloped] {keywords} (Key);
                            \draw[arrow] (Title) -- node[labelstyle, sloped] {belongs to} (Field);

                            % --- Title to Concepts ---
                            \draw[arrow] (Title) -- node[labelstyle, sloped, pos=0.5] {solves} (Prob);
                            \draw[arrow] (Title) -- node[labelstyle, sloped, pos=0.5] {adopts} (Method);
                            \draw[arrow] (Title) -- node[labelstyle, sloped, pos=0.5] {uses} (Metric);
                            \draw[arrow] (Title) -- node[labelstyle, sloped, pos=0.25] {experiments on} (Dataset);
                            \draw[arrow] (Title) -- node[labelstyle, sloped, pos=0.3] {proposes} (Model);
                            \draw[arrow] (Title) -- node[labelstyle, sloped, pos=0.4] {has} (Results);

                            % --- Inter-concept Connections ---

                            % Method relations
                            \draw[arrow] (Method) -- node[labelstyle, pos=0.5] {solves} (Prob);
                            \draw[arrow] (Method) -- node[labelstyle, sloped, pos=0.5] {works on} (Task);
                            \draw[arrow] (Method) -- node[labelstyle, sloped, pos=0.6] {proposes} (Model);
                            \draw[arrow] (Method) -- node[labelstyle, sloped, pos=0.25] {innovates} (Innov);
                            \draw[arrow] (Method) -- node[labelstyle, sloped, pos=0.25] {has} (Results);

                            % Task relations
                            \draw[arrow] (Task) -- node[labelstyle, sloped, pos=0.25] {faces} (Prob);
                            \draw[arrow] (Task) -- node[labelstyle, sloped, pos=0.55] {uses} (Metric);
                            \draw[arrow] (Task) -- node[labelstyle, sloped, pos=0.5] {works\\for} (Model);
                            \draw[arrow] (Task) -- node[labelstyle, sloped, pos=0.5] {experiments on} (Dataset);

                            % Model relations
                            \draw[arrow] (Model) -- node[labelstyle, pos=0.50] {has} (Results);
                            \draw[arrow] (Model) -- node[labelstyle, sloped, pos=0.7] {solves} (Prob);
                            \draw[arrow] (Model) -- node[labelstyle, sloped, pos=0.25] {uses} (Dataset);

                            % Results relations
                            \draw[arrow] (Results) -- node[labelstyle, sloped, pos=0.25] {uses} (Metric);

                            % Innovation relations
                            \draw[arrow] (Innov) -- node[labelstyle, sloped, pos=0.5] {uses} (Dataset);

                        \end{tikzpicture}
                        \caption{Skema ontologi (\textit{Ontology Design}) \textit{Knowledge Graph} untuk domain riset akademik}
                        \label{fig:ontology-design}
                    \end{figure}

                    \hspace*{0.75cm}Skema ontologi dirancang sebagai \textit{blueprint} \textit{Knowledge Graph} sebelum proses ekstraksi dimulai (\ref{fig:ontology-design}). Implementasi difokuskan pada delapan relasi utama: \textit{publishes}, \textit{cites}, \textit{has\_author}, \textit{has\_topic}, \textit{uses\_method}, \textit{has\_affiliation}, \textit{belongs\_to\_domain}, dan \textit{collaborates\_with}. Konsistensi skema ini membatasi ruang pencarian model ekstraksi agar tetap relevan dengan domain Infokom UNESA.

              \item \textbf{\textit{Parallel Processing Pipeline}}

                    \hspace*{0.75cm}Data akademik diproses melalui dua jalur paralel yang dioptimalkan untuk karakteristik data berbeda. Kedua jalur berjalan independen namun menghasilkan keluaran yang saling melengkapi (\ref{fig:detail-kg-construction}).

                    \begin{enumerate}[label=\textbf{\alph*.}, leftmargin=0.75cm, itemsep=5pt]

                        \item \textbf{Jalur Data Terstruktur: \textit{Entity Resolution}}

                              \hspace*{0.75cm}Data terstruktur (profil dosen dan afiliasi) diproses melalui \textit{Entity Resolution} untuk menyatukan referensi duplikat. Variasi penulisan nama (misalnya ``Aditya Prapanca'' vs. ``A. Prapanca'') diselesaikan untuk menghasilkan satu \textit{canonical node} per dosen, menjamin akumulasi publikasi yang akurat pada satu entitas.

                        \item \textbf{Jalur Data Tidak Terstruktur: \textit{NLP-Based Extraction}}

                              \hspace*{0.75cm}Tahap awal \textbf{\textit{Text Processing} \& \textit{Coreference Resolution}} membersihkan abstrak dari \textit{noise} dan mengganti kata ganti (``metode ini'', ``itu'') dengan entitas referensinya untuk mempertahankan konteks penuh.

                              \hspace*{0.75cm}Selanjutnya, \textbf{Ekstraksi Entitas dan Relasi} dilakukan menggunakan \textit{GLiNER} \cite{zaratiana2023gliner} untuk \textit{Named Entity Recognition} secara \textit{zero-shot} dan \textit{GLiRel} \cite{boylan2025glirelgeneralistmodel} untuk ekstraksi relasi. Kombinasi ini dipilih karena tidak memerlukan dataset pelatihan teranotasi. Keluaran berupa \textit{Lexical Graph} berisi triplet kandidat.

                    \end{enumerate}

              \item \textbf{\textit{Entity Linking}}

                    \hspace*{0.75cm}Tahap ini menjembatani kedua jalur dengan memetakan entitas hasil ekstraksi teks (nama dosen, topik riset) ke entitas kanonikal dari jalur terstruktur. Misalnya, penyebutan ``Deep Learning'' dalam abstrak dihubungkan ke node topik yang sudah ada, bukan membentuk node duplikat. Hasilnya adalah \textit{Scholarly Knowledge Graph} utuh yang mengintegrasikan entitas terverifikasi dengan relasi dari teks tanpa duplikasi.

              \item \textbf{\textit{Dual Indexing}}

                    \hspace*{0.75cm}Untuk mendukung \textit{Hybrid Retrieval}, hasil konstruksi \textit{KG} disimpan dalam dua format berikut:

                    \begin{enumerate}[label=\textbf{\alph*.}, leftmargin=0.75cm, itemsep=5pt]
                        \item \textbf{\textit{Vector Indexing}}

                              \hspace*{0.75cm}Teks abstrak dipotong menjadi segmen 2.024 karakter dengan \textit{overlap} 50 karakter, dikonversi ke vektor \textit{embedding} menggunakan \texttt{all-MiniLM-L6-v2} (384 dimensi), dan disimpan di \textit{Weaviate} untuk pencarian kemiripan semantik.

                        \item \textbf{\textit{Graph Indexing}}

                              \hspace*{0.75cm}Triplet entitas dan relasi dimuat ke basis data graf \textit{Neo4j} melalui \textit{Cypher queries}, memfasilitasi pencarian berbasis struktur relasional dan \textit{graph traversal}.
                    \end{enumerate}

          \end{enumerate}

    \item \textit{\textbf{Hybrid Retrieval pada GraphRAG Model}}

          \hspace*{0.75cm}Berbeda dengan \textit{RAG} konvensional yang hanya mengandalkan pencarian vektor, pendekatan ini memperkuat \textit{retrieval} dengan struktur \textit{Knowledge Graph} (\ref{fig:detail-hybrid-rag}). Kueri pengguna diproses melalui serangkaian tahap optimasi sebelum pengambilan data.

          \begin{figure}[htbp]
              \centering
              \begin{tikzpicture}[scale=1, transform shape]
                  % Outer box — widened so Local/Global nodes (2.0cm) fit with margin
                  \node [subbox, fit={(-0.5,0.3) (5.3,-9.3)}, label={[font=\scriptsize\bfseries]above:Detail Hybrid Graph-RAG}] {};

                  \node (subD3) [offpage] at (2.4, -0.1) {D};
                  \node (qin) [data] at (2.4, -1.3) {User Query};
                  \node (clue) [process, text width=2.2cm] at (2.4, -2.5) {Semantic Clue\\Extraction};
                  \node (kwd) [process, text width=2.2cm] at (2.4, -3.8) {Keyword\\Decomposition};
                  % Local/Global — override minimum width to 2.0cm (default 2.8cm overflows box)
                  \node (local) [process, text width=1.6cm, minimum width=2.0cm] at (1.0, -5.2) {Local\\Retrieval};
                  \node (globalret) [process, text width=1.6cm, minimum width=2.0cm] at (3.8, -5.2) {Global\\Retrieval};
                  \node (integrate) [process, text width=2.2cm] at (2.4, -6.5) {Context\\Fusion};
                  \node (llm) [process, text width=2.4cm] at (2.4, -7.7) {LLM Generation};
                  \node (response) [data] at (2.4, -8.9) {Response};

                  \draw [arrow] (subD3) -- (qin);
                  \draw [arrow] (qin) -- (clue);
                  \draw [arrow] (clue) -- (kwd);
                  \draw [arrow] (kwd) -- (local);
                  \draw [arrow] (kwd) -- (globalret);
                  \draw [arrow] (local) -- (integrate);
                  \draw [arrow] (globalret) -- (integrate);
                  \draw [arrow] (integrate) -- (llm);
                  \draw [arrow] (llm) -- (response);
              \end{tikzpicture}
              \caption{Detail proses \textit{Hybrid Retrieval}}
              \label{fig:detail-hybrid-rag}
          \end{figure}

          \begin{enumerate}[label=\textbf{\alph*.}, leftmargin=*, itemsep=5pt]

              \item \textbf{Pengambilan Petunjuk Semantik (\textit{Clues}) pada \textit{Vector Database}}

                    \hspace*{0.75cm}Tahap ini mengidentifikasi entitas kunci dari pertanyaan pengguna via pencarian vektor, lalu memetakannya ke terminologi ontologi graf (misalnya ``AI'' dipetakan ke ``Kecerdasan Buatan'').

                    \hspace*{0.75cm}Pertanyaan pengguna $Q$ dipetakan ke ruang vektor menggunakan fungsi \textit{embedding}, kemudian diukur kedekatannya dengan \textit{keyword} korpus via \textit{cosine similarity}. Seleksi kandidat \textit{clues} didefinisikan sebagai:

                    \begin{equation}
                        K^{\text{clues}} = \{k \in K^{\text{content}} \mid \text{sim}(f_{\text{embed}}(Q), f_{\text{embed}}(k)) \geq \tau\}
                    \end{equation}

                    \hspace*{0.75cm}Setelah \textit{clues} diekstraksi, model bahasa besar (\textit{LLM}) mengubahnya menjadi dua tingkat kata kunci yang saling melengkapi:

                    \begin{equation}
                        (\mathcal{K}_h, \mathcal{K}_l) = \text{LLM}_{keyword}(Q, K^{\text{clues}})
                    \end{equation}

                    Kata kunci \textit{High-level} $\mathcal{K}_h$ menangkap konsep abstrak (misalnya ``Deep Learning'', ``Sequence Processing''), sementara \textit{Low-level} $\mathcal{K}_l$ menangkap detail teknis spesifik (misalnya ``Attention Mechanism'', ``Encoder-Decoder Models''). Strategi dua tingkat ini memungkinkan sistem menjawab pertanyaan holistik ($\mathcal{K}_h$) maupun presisi tinggi ($\mathcal{K}_l$).
                    \begin{figure}[htbp]
                        \centering
                        \begin{tikzpicture}[
                                >=Stealth,
                                scale=1, transform shape,
                                % Vivid color scheme matching reference image
                                userbox/.style={rectangle, draw=blue!80!black, fill=blue!70!black, text=white, minimum width=2.4cm, minimum height=1.4cm, align=center, font=\small\bfseries, rounded corners=6pt, line width=0.7pt},
                                cluesbox/.style={rectangle, draw=orange!80!black, fill=orange!30, minimum width=2.4cm, minimum height=1.5cm, align=center, font=\small, rounded corners=6pt, line width=0.7pt},
                                lowlevelbox/.style={rectangle, draw=green!60!black, fill=green!30, minimum width=2.4cm, minimum height=1.5cm, align=center, font=\small, rounded corners=6pt, line width=0.7pt},
                                highlevelbox/.style={rectangle, draw=red!60!black, fill=red!20, minimum width=2.4cm, minimum height=1.5cm, align=center, font=\small, rounded corners=6pt, line width=0.7pt},
                                arrow/.style={->, >=Stealth, line width=0.8pt, blue!70!cyan},
                                label/.style={font=\scriptsize, align=left}
                            ]

                            % === LEFT SIDE (Input) ===
                            % User Query Box (soft blue) -> shifted in to -3.8
                            \node[userbox] (user) at (-3.8, 1.5) {\textbf{User:} How does\\transformer\\architecture work?};

                            % Clues Box (soft orange/peach) -> shifted in to -3.8
                            \node[cluesbox] (clues) at (-3.8, -1.5) {Model architecture,\\Transformer,\\Transformer\\architecture...};

                            % === CENTER: LLM Robot (at x=0) ===
                            % Robot head (main rectangle)
                            \draw[fill=black!80, draw=black, line width=0.6pt, rounded corners=3pt] (-0.4, -0.1) rectangle (0.4, 0.7);
                            % Antenna
                            \draw[black!80, line width=1.2pt] (0, 0.7) -- (0, 0.95);
                            \fill[black!80] (0, 0.95) circle (0.1);
                            % Eyes
                            \draw[fill=white, draw=black, line width=0.4pt] (-0.22, 0.38) circle (0.14);
                            \draw[fill=white, draw=black, line width=0.4pt] (0.12, 0.38) circle (0.14);
                            \fill[black] (-0.22, 0.38) circle (0.06);
                            \fill[black] (0.12, 0.38) circle (0.06);
                            % Ears
                            \draw[fill=black!70, draw=black, line width=0.4pt] (-0.5, 0.15) rectangle (-0.4, 0.55);
                            \draw[fill=black!70, draw=black, line width=0.4pt] (0.4, 0.15) rectangle (0.5, 0.55);
                            % Body
                            \draw[fill=black!80, draw=black, line width=0.5pt, rounded corners=2pt] (-0.3, -0.55) rectangle (0.3, -0.1);
                            % Body lines
                            \draw[gray!45, line width=0.4pt] (-0.2, -0.25) -- (0.2, -0.25);
                            \draw[gray!45, line width=0.4pt] (-0.2, -0.4) -- (0.2, -0.4);

                            % LLM label
                            \node[font=\normalsize\bfseries, text=blue!70!black] at (0, -0.9) {LLM};

                            % === RIGHT SIDE (Output) ===
                            % Low-level keywords box (soft green) - top right -> shifted in to 3.8
                            \node[lowlevelbox] (lowlevel) at (3.8, 1.5) {Attention mechanism,\\Encoder-decoder\\models, Language\\modeling...};

                            % High-level keywords box (soft pink) - bottom right -> shifted in to 3.8
                            \node[highlevelbox] (highlevel) at (3.8, -1.5) {Transformer\\architecture, Deep\\learning, Sequence\\processing...};

                            % === ARROWS ===

                            % ① User -> Clues (down arrow)
                            \draw[arrow] (user.south) -- (clues.north);

                            % User -> LLM (diagonal down-right)
                            \draw[arrow] (user.east) -- (-0.6, 0.5);

                            % Clues -> LLM (diagonal up-right)
                            \draw[arrow] (clues.east) -- (-0.6, -0.3);

                            % LLM -> Low-level (diagonal up-right)
                            \draw[arrow] (0.6, 0.5) -- node[midway, above, sloped, font=\scriptsize\itshape, text=red!50!gray, yshift=2pt] {Low-level} (lowlevel.west);

                            % LLM -> High-level (diagonal down-right)
                            \draw[arrow] (0.6, -0.3) -- node[midway, below, sloped, font=\scriptsize\itshape, text=red!50!gray, yshift=-2pt] {High-level} (highlevel.west);

                            % === LABELS ===

                            % ① Retrieve Content Keywords (left of arrow) -> adjusted x
                            \node[label, anchor=east] at (-4.1, 0) {\textcircled{\scriptsize 1} \textbf{Retrieve Content}\\Keywords by Similarity};

                            % ② Input to LLM (centered above LLM)
                            \node[label, anchor=south] at (0, 1.15) {\textcircled{\scriptsize 2} \textbf{Input to LLM}};

                            % ③ Generate Hierarchical Keywords (right side) -> adjusted x
                            \node[label, anchor=west] at (1.6, 0) {\textcircled{\scriptsize 3} \textbf{Generate Hierarchical}\\~~~~\textbf{Keywords}};

                        \end{tikzpicture}
                        \caption{Ilustrasi Ekstraksi Kata Kunci Berdasarkan Petunjuk (\textit{clues}) \citep{chen2025acadrag}}
                        \label{fig:contoh_ekstraksi_kata_kunci}
                    \end{figure}

                    \hspace*{0.75cm}\ref{fig:contoh_ekstraksi_kata_kunci} mengilustrasikan proses ini. \textit{LLM} berfungsi sebagai \textit{reasoning engine} yang mensintesis pertanyaan asli dengan petunjuk semantik, memecah kompleksitas kueri menjadi komponen pencarian yang relevan secara leksikal dan kontekstual. Pendekatan hierarkis ini meningkatkan \textit{recall} sekaligus menjaga \textit{precision}.

              \item \textbf{Ekstraksi Informasi Lokal Berbasis Subgraf (\textit{Local Retrieval})}
                    % Diagram Tahapan Pembentukan Subgraf (TikZ) - Academic Style
                    \begin{figure}[htbp]
                        \centering
                        \begin{tikzpicture}[
                                >=Stealth,
                                scale=1, transform shape,
                                % Vivid Green/Blue color scheme matching reference
                                graphnode/.style={circle, draw=blue!40!black, fill=green!60!black, minimum size=0.45cm, line width=0.7pt},
                                graphnodelight/.style={circle, draw=blue!40!black, fill=green!15, minimum size=0.45cm, line width=0.7pt},
                                edge/.style={line width=1.5pt, blue!60!cyan},
                                prunededge/.style={line width=1.5pt, red!80!black},
                                newedge/.style={line width=1.8pt, red!80!black},
                                arrow/.style={->, line width=0.7pt, black!80},
                                dashedoval/.style={draw=orange!90!black, densely dashed, line width=1.0pt}
                            ]

                            % ========== STAGE 1: Node Matching (3 nodes) ==========
                            \begin{scope}[shift={(0,0)}]
                                % Three separate nodes (tidak terhubung)
                                \node[graphnode] (n1a) at (0, 0.6) {};
                                \node[graphnode] (n1b) at (0.4, -0.2) {};
                                \node[graphnode] (n1c) at (-0.3, -0.5) {};

                                % Label
                                \node[font=\scriptsize\itshape] at (0, -1.2) {Node Matching};
                            \end{scope}

                            % Arrow to next stage
                            \draw[arrow] (1.0, 0) -- (1.5, 0);

                            % ========== STAGE 2: Subgraph Construction (5 nodes pentagon) ==========
                            \begin{scope}[shift={(2.8, 0)}]
                                % Pentagon graph - 5 nodes with "Potential Missing Node" pointing here
                                \node[graphnodelight] (n2a) at (0, 0.8) {};      % top (light) - target of potential missing node
                                \node[graphnode] (n2b) at (0.65, 0.25) {};       % right upper (dark)
                                \node[graphnodelight] (n2c) at (0.45, -0.5) {};  % right lower (light)
                                \node[graphnode] (n2d) at (-0.35, -0.6) {};      % bottom left (dark)
                                \node[graphnode] (n2e) at (-0.6, 0.1) {};        % left (dark)

                                % Edges (blue) - pentagon + diagonal
                                \draw[edge] (n2a) -- (n2b);
                                \draw[edge] (n2b) -- (n2c);
                                \draw[edge] (n2c) -- (n2d);
                                \draw[edge] (n2d) -- (n2e);
                                \draw[edge] (n2e) -- (n2a);

                                % Dashed oval "Potential Missing Node" - pointing to top node of subgraph
                                \draw[dashedoval] (-1.5, 1.6) ellipse (0.75cm and 0.4cm);
                                \node[font=\tiny\itshape, align=center, text=orange!70!black] at (-1.5, 1.6) {Potential\\Missing Node};

                                % Curved dashed arrow from oval to top node of subgraph
                                \draw[orange!70!black, densely dashed, line width=0.9pt, ->]
                                (-1.5, 1.15) to[out=-90, in=90] (n2e.north);

                                % Label
                                \node[font=\scriptsize\itshape, align=center] at (0, -1.2) {Subgraph\\Construction};
                            \end{scope}

                            % Arrow to next stage
                            \draw[arrow] (4.2, 0) -- (4.7, 0);

                            % ========== STAGE 3: Pruning (5 nodes, 1 isolated node with slash) ==========
                            \begin{scope}[shift={(6.0, 0)}]
                                % Pentagon - 5 nodes, one is isolated (not connected)
                                \node[graphnode] (n3a) at (-0.15, 0.8) {};       % top left (dark)
                                \node[graphnodelight] (n3b) at (0.55, 0.55) {};  % top right (light) - ISOLATED, will be pruned
                                \node[graphnode] (n3c) at (0.65, -0.15) {};      % right (dark)
                                \node[graphnode] (n3d) at (0, -0.7) {};          % bottom (dark)
                                \node[graphnodelight] (n3e) at (-0.6, -0.1) {};  % left (light)

                                % Blue edges - connecting all EXCEPT n3b (isolated node)
                                \draw[edge] (n3a) -- (n3e);
                                \draw[edge] (n3e) -- (n3d);
                                \draw[edge] (n3d) -- (n3c);

                                % Red slash mark ON TOP of the isolated node (n3b) - marking it for pruning
                                \draw[red!70!black, line width=2.5pt] (0.38, 0.72) -- (0.72, 0.38);

                                % Label
                                \node[font=\scriptsize\itshape] at (0, -1.2) {Pruning};
                            \end{scope}

                            % Arrow to next stage
                            \draw[arrow] (7.4, 0) -- (7.9, 0);

                            % ========== STAGE 4: Edge Expansion (5 nodes, 1 disconnected) ==========
                            \begin{scope}[shift={(9.2, 0)}]
                                % 5 nodes - one is disconnected (isolated)
                                \node[graphnode] (n4a) at (0, 0.75) {};          % top (dark)
                                \node[graphnodelight] (n4b) at (0.6, 0.15) {};   % right upper (light) - DISCONNECTED
                                \node[graphnode] (n4c) at (0.4, -0.55) {};       % right lower (dark)
                                \node[graphnodelight] (n4d) at (-0.4, -0.55) {}; % left lower (light)
                                \node[graphnode] (n4e) at (-0.6, 0.15) {};       % left upper (dark)

                                % Blue edges - NOT connecting n4b (disconnected node)
                                \draw[edge] (n4a) -- (n4e);
                                \draw[edge] (n4e) -- (n4d);
                                \draw[edge] (n4d) -- (n4c);

                                % Red expansion edges (outward from each connected node)
                                % From top node
                                \draw[newedge] (n4a) -- ++(0.22, 0.32);
                                \draw[newedge] (n4a) -- ++(-0.18, 0.35);
                                % From right upper (disconnected - still has expansion edges)
                                \draw[newedge] (n4b) -- ++(0.32, 0.12);
                                \draw[newedge] (n4b) -- ++(0.28, -0.18);
                                % From right lower
                                \draw[newedge] (n4c) -- ++(0.28, -0.22);
                                \draw[newedge] (n4c) -- ++(0.06, -0.38);
                                % From left lower
                                \draw[newedge] (n4d) -- ++(-0.28, -0.22);
                                \draw[newedge] (n4d) -- ++(-0.06, -0.38);
                                % From left upper
                                \draw[newedge] (n4e) -- ++(-0.32, 0.12);
                                \draw[newedge] (n4e) -- ++(-0.28, -0.18);

                                % Label
                                \node[font=\scriptsize\itshape, align=center] at (0, -1.2) {Edge\\Expansion};
                            \end{scope}

                        \end{tikzpicture}
                        \caption{Tahapan Pembentukan Subgraf}
                        \label{fig:pembentukan-subgraf}
                    \end{figure}

                    \hspace*{0.75cm}Tahap ini menjawab pertanyaan yang membutuhkan ketepatan tinggi terhadap entitas tertentu. Sistem melakukan \textit{anchor node matching} untuk menemukan entitas relevan, lalu menelusuri graf sejauh 1--2 hop untuk mengambil konteks tetangga. Metode ekstraksi subgraf diarahkan oleh kata kunci \textit{low-level} $\mathcal{K}_l$ dengan tahapan berikut:

                    \begin{enumerate}[label=\textbf{\arabic*.}, leftmargin=*, itemsep=6pt]

                        \item \textbf{\textit{Node Matching} (Pencocokan Entitas):} \textit{Graph nodes} dicocokkan dengan kata kunci \textit{low level} ($\mathcal{K}_l$) menggunakan metrik \textit{vector similarity} yang memprioritaskan \textit{nodes} dengan nilai \textit{cosine similarity} tertinggi terhadap \textit{keyword embeddings}.

                        \item \textbf{\textit{Subgraph Construction}:} Simpul-simpul jangkar dihubungkan menggunakan \textit{shortest path algorithms} untuk menemukan rute paling efisien antar entitas terkait.

                        \item \textbf{\textit{Pruning}:} \textit{Node} atau \textit{edge} dengan relevansi rendah dieliminasi agar subgraf tetap padat dan bermakna.

                        \item \textbf{\textit{Text Unit Extraction}:} \textit{Textual chunks} yang terasosiasi dengan \textit{retained edges} dalam lingkup \textit{one-hop ego-graph} diperingkat berdasarkan frekuensi kemunculannya. Segmen teratas diprioritaskan untuk \textit{retrieval}.

                        \item \textbf{\textit{Edge Expansion}:} Subgraf diperluas dengan menggabungkan \textit{one-hop neighbors} dari \textit{matched nodes} beserta \textit{edges} yang memenuhi kriteria relevansi.

                        \item \textbf{\textit{Data Truncation}:} Data yang diekstraksi (\textit{nodes}, \textit{edges}, \textit{text units}) dipangkas hingga total token memenuhi ambang batas komputasi.

                    \end{enumerate}

                    \hspace*{0.75cm}Proses ini menghasilkan konteks lokal terstruktur (\ref{fig:pembentukan-subgraf}). Hasil akhir dinotasikan sebagai tiga komponen: entitas ($\mathcal{E}$), relasi ($\mathcal{R}$), dan unit teks ($\mathcal{T}$):

                    \begin{equation}
                        \mathcal{L}_{\text{local}} = (\mathcal{E}_{\text{local}}, \mathcal{R}_{\text{local}}, \mathcal{T}_{\text{local}})
                        = \text{pipeline}_{\text{local}}(\mathcal{K}_l, G).
                    \end{equation}

                    \noindent
                    di mana $\mathcal{L}_{\text{local}}$ merupakan fungsi dari kata kunci spesifik dan struktur graf, yang direduksi menjadi subset informasi paling signifikan untuk menjawab kueri.


              \item \textbf{Ekstraksi Jejaring Global (\textit{Global Edge Network Retrieval})}

                    \hspace*{0.75cm}Jalur ini menjawab pertanyaan eksploratif yang membutuhkan pemahaman tematik luas (\textit{sensemaking}). Sistem menggunakan \textit{High-level Keywords} $\mathcal{K}_h$ untuk mengekstraksi \textit{Global Edge Network} berbasis vektor, tanpa batasan jarak topologi antar \textit{nodes}. Tahapannya meliputi:

                    \begin{enumerate}[label=\textbf{\arabic*.}, leftmargin=*, itemsep=6pt]

                        \item \textbf{\textit{Edge Retrieval}:} Basis data vektor dikueri untuk mengidentifikasi \textit{edges} yang representasi vektornya berkorelasi dengan kata kunci tingkat tinggi, dengan memprioritaskan \textit{edges} yang merangkum hubungan konseptual makro (seperti kerangka teori atau tren lintas domain).

                        \item \textbf{\textit{Nodes Retrieval}:} \textit{Nodes} yang terkait dengan \textit{edges} terpilih diekstraksi, kemudian diurutkan berdasarkan derajat keterhubungannya (\textit{degree}).

                        \item \textbf{\textit{Edge Data Processing}:} Metadata \textit{edges} diekstraksi, meliputi \textit{degree} (jumlah koneksi \textit{source/target node}) dan \textit{weight} (kekuatan hubungan semantik) untuk mengkuantifikasi signifikansinya.

                        \item \textbf{\textit{Integration and Ranking}:} \textit{Edges} diperingkat menggunakan metrik komposit yang menggabungkan skor relevansi dan bobot koneksi, memprioritaskan tautan konseptual dominan.

                        \item \textbf{\textit{Text Unit Extraction}:} \textit{Text chunks} terasosiasi dengan \textit{edges} diperingkat berdasarkan kemiripan dengan kata kunci \textit{global} dan diambil sebagai konteks pendukung.

                        \item \textbf{\textit{Data Truncation}:} Data dipangkas hingga total token memenuhi ambang batas komputasi.

                    \end{enumerate}

                    \noindent
                    Mekanisme ini menghasilkan konteks global yang kaya akan informasi hubungan, yang diformulasikan dalam persamaan:
                    \begin{equation}
                        \mathcal{L}_{\text{global}} = (\mathcal{E}_{\text{global}}, \mathcal{R}_{\text{global}}, \mathcal{T}_{\text{global}})
                        = \text{pipeline}_{\text{global}}(\mathcal{K}_h, G).
                    \end{equation}
                    \noindent
                    Sinergi antara presisi tinggi \textit{Local Retrieval} dan cakupan luas \textit{Global Retrieval} menjadi keunggulan utama arsitektur \textit{GraphRAG} dalam menangani kompleksitas literatur ilmiah.

              \item \textbf{\textit{LLM Generation} \& \textit{Response}}
                    \hspace*{0.75cm}Tahap ini merupakan fase akhir, di mana informasi yang telah diekstraksi dari \textit{local} dan \textit{global} disintesis. Konteks yang terintegrasi ini kemudian disatukan dengan kueri asli pengguna untuk membentuk sebuah \textit{augmented prompt}. \textit{Augmented prompt} ini kemudian dikirimkan ke \textit{LLM} untuk menghasilkan jawaban akhir.
          \end{enumerate}

    \item \textbf{Evaluasi Sistem}
          \label{subsec:system-evaluation}

          \hspace*{0.75cm}Penelitian ini mengevaluasi kinerja sistem menggunakan kerangka \textit{RAGAS (Retrieval-Augmented Generation Assessment System)} yang menyediakan metrik komprehensif untuk menilai dua dimensi utama, yaitu: kemampuan pengambilan informasi relevan (\textit{retrieval quality}) dan kualitas jawaban yang dihasilkan (\textit{generation quality}).

          \hspace*{0.75cm}Skenario eksperimen dirancang untuk membandingkan \textit{GraphRAG} dan \textit{Vector RAG} menggunakan 40 pertanyaan akademik yang mencakup aspek faktual, relasional, dan sintesis, yang masing-masing dilengkapi dengan \textit{ground truth} manual. Lima metrik utama digunakan sebagai tolok ukur:

          \begin{enumerate}[label=\arabic*., leftmargin=0.75cm, itemsep=3pt]

              \item \textbf{\textit{Context Precision}}: Relevansi dokumen terhadap pertanyaan
                    \begin{equation}
                        \textit{Context Precision} = \frac{\sum_{k=1}^{K} \textit{Precision}(k) \cdot v_k}{\text{Total item relevan}}
                    \end{equation}

              \item \textbf{\textit{Context Recall}}: Rasio \textit{ground truth} yang didukung konteks
                    \begin{equation}
                        \textit{Context Recall} = \frac{\text{Kalimat GT didukung}}{\text{Total kalimat GT}}
                    \end{equation}

              \item \textbf{\textit{Answer Correctness}}: Kecocokan jawaban dengan \textit{ground truth}
                    \begin{equation}
                        \textit{Answer Correctness} = \frac{\text{TP}}{\text{TP} + 0.5 \cdot \text{FP} + \text{FN}}
                    \end{equation}

              \item \textbf{\textit{Faithfulness}}: Rasio klaim terverifikasi (kontrol halusinasi)
                    \begin{equation}
                        \textit{Faithfulness} = \frac{|V|}{|S|}
                    \end{equation}

              \item \textbf{\textit{Answer Relevancy}}: Keselarasan jawaban dengan pertanyaan
                    \begin{equation}
                        \textit{Answer Relevancy} = \frac{1}{n} \sum_{i=1}^{n} \cos(E_{q_i}, E_q)
                    \end{equation}

          \end{enumerate}
          Penelitian menetapkan target keberhasilan sistem berdasarkan \textit{threshold} berikut:
          \begin{itemize}[leftmargin=0.75cm]
              \item \textit{Context Precision} $> 0.85$ (relevansi tinggi)
              \item \textit{Context Recall} $> 0.80$ (cakupan informasi luas)
              \item \textit{Faithfulness} $> 0.90$ (minimal halusinasi)
              \item \textit{Answer Correctness} $> 0.75$ (keakuratan jawaban)
          \end{itemize}

          Sistem dianggap berhasil apabila memenuhi standar seluruh metrik tersebut. Untuk memvalidasi efektivitas, dilakukan perbandingan \textit{benchmarking} antara \textit{Baseline} (\textit{Standard Vector RAG}) dengan sistem \textit{Hybrid GraphRAG} yang diusulkan.


          % 6. GUI Implementation
    \item \textbf{Implementasi GUI (\textit{GUI Implementation})}

          \hspace*{0.75cm}Antarmuka pengguna (\textit{Graphical User Interface}) dikembangkan menggunakan \textit{Streamlit} untuk memfasilitasi interaksi pengguna dengan sistem. Antarmuka ini menampilkan jawaban, referensi sumber, dan visualisasi subgraf secara \textit{real-time}, sehingga memberikan transparansi terhadap proses perolehan jawaban.
\end{enumerate}

\section{JADWAL PENELITIAN}

Jadwal penelitian disusun dalam \ref{tab:jadwal-penelitian}, mencakup seluruh tahapan dari persiapan hingga penyelesaian skripsi selama enam bulan, namun tetap fleksibel menyesuaikan perkembangan progres.

\begin{table}[htbp]
    \centering
    \caption{Pelaksanaan Penelitian}
    \label{tab:jadwal-penelitian}
    \renewcommand{\arraystretch}{1.3}
    \small
    \begin{tabular}{@{}|l|c|c|c|c|c|c|@{}}
        \hline
        \textbf{Kegiatan}       & \textbf{Januari}    & \textbf{Februari}   & \textbf{Maret}      & \textbf{April}      & \textbf{Mei}        & \textbf{Juni}       \\
        \hline
        Studi Literatur         & \cellcolor{gray!40} & \cellcolor{gray!40} &                     &                     &                     &                     \\
        \hline
        Pengumpulan Data        &                     & \cellcolor{gray!40} & \cellcolor{gray!40} &                     &                     &                     \\
        \hline
        Persiapan Data          &                     &                     & \cellcolor{gray!40} & \cellcolor{gray!40} &                     &                     \\
        \hline
        Konstruksi KG           &                     &                     &                     & \cellcolor{gray!40} & \cellcolor{gray!40} &                     \\
        \hline
        Implementasi Hybrid RAG &                     &                     &                     &                     & \cellcolor{gray!40} & \cellcolor{gray!40} \\
        \hline
        Evaluasi Sistem         &                     &                     &                     &                     &                     & \cellcolor{gray!40} \\
        \hline
        Penulisan Skripsi       &                     & \cellcolor{gray!40} & \cellcolor{gray!40} & \cellcolor{gray!40} & \cellcolor{gray!40} & \cellcolor{gray!40} \\
        \hline
    \end{tabular}
\end{table}
